\documentclass[tikz,border=4pt]{standalone}
\usepackage{ctex}
\usepackage{amsmath}
\usepackage{pgfplots}
\pgfplotsset{compat=1.18}

% p10-03b: 二次函数性质图——a<0，对称轴左侧递增，右侧递减
% y = -(x-1)^2 + 3，对称轴x=1，顶点(1,3)
\begin{document}
\begin{tikzpicture}
\begin{axis}[
  width=8cm, height=7.5cm,
  axis lines=center,
  xlabel={$x$}, ylabel={$y$},
  every axis x label/.style={at={(axis cs:4.5,0)}, anchor=west},
  every axis y label/.style={at={(axis cs:0,4.5)}, anchor=south},
  xmin=-2.0, xmax=4.5,
  ymin=-2.0, ymax=4.5,
  xtick={-2,-1,1,2,3,4},
  ytick={-2,-1,1,2,3,4},
  tick style={color=black},
  ticklabel style={font=\small},
  clip=false,
]
  % y = -(x-1)^2 + 3
  \addplot[blue, thick, domain=-1.7:3.7, samples=100] {-(x-1)^2 + 3};
  \node[right, blue] at (axis cs:3.0, 0.2) {$y=-(x-1)^2+3$};

  % 对称轴 x=1（虚线）
  \draw[gray, dashed, thick] (axis cs:1,-2.0) -- (axis cs:1,4.5);
  \node[above, gray, font=\small] at (axis cs:1,4.5) {$x=1$};

  % 顶点 (1,3)
  \addplot[only marks, mark=*, mark size=3pt, red] coordinates {(1,3)};
  \node[above left, red] at (axis cs:1,3) {$(1,\,3)$};

  % 递增区间箭头（x<1，从左向右）
  \draw[->, blue!70, very thick] (axis cs:-1.0, -1.0) -- (axis cs:0.9, 3.0) node[midway, left, blue, font=\small] {递增};

  % 递减区间箭头（x>1，从左向右向下）
  \draw[->, blue!70, very thick] (axis cs:1.1, 3.0) -- (axis cs:3.0, -1.0) node[midway, right, blue, font=\small] {递减};

  % 区间标注
  \node[below, font=\small, blue] at (axis cs:-0.5, -0.15) {$x<1$};
  \node[below, font=\small, blue] at (axis cs:2.5, -0.15) {$x>1$};
\end{axis}
\end{tikzpicture}
\end{document}
